\documentclass[11pt, a4paper]{article}

% --- EVRENSEL ÖN BİLGİ VE PAKET TANIMLARI ---
\usepackage[a4paper, top=2.5cm, bottom=2.5cm, left=2cm, right=2cm]{geometry}
\usepackage{fontspec}

\usepackage[turkish, bidi=basic, provide=*]{babel}
\babelprovide[import, onchar=ids fonts]{turkish}
\babelprovide[import, onchar=ids fonts]{english}

\babelfont{rm}{Noto Sans}
\babelfont[turkish]{rm}{Noto Sans}

\usepackage{enumitem}
\setlist[itemize]{label=-}

\usepackage{amsmath,amssymb,amsthm,mathtools,bm,mathrsfs}
\usepackage{physics}
\usepackage{booktabs}

\usepackage{hyperref}
\hypersetup{
    colorlinks=true,
    linkcolor=blue,
    citecolor=blue,
    urlcolor=blue
}

\title{\textbf{Nefs-i Müdrike Mimarisi: Sanal Sayıların ($i = \sqrt{-1}$) İlgası ve 41 Melekenin Hakiki Reel Kuantum Operatörleri ile Geometrik Cebir Formülasyonu}}
\author{\textbf{Reel Geometrik Cebir ($Cl_{3,0}$), Reel Ortogonal Gruplar ($SO(N)$), Reel Hartley Spektral Süzgeçleri ve $p$-Adic Rasyonel İntaç Nazariyesi}}
\date{\today}

\begin{document}

\maketitle

\begin{abstract}
Bu risale; kuantum hesaplamaya sızdırılan soyut ve hayali sanal sayıları ($i = \sqrt{-1}$, karmaşık karmaşasını) kâmilen ilga ederek, Nefs-i Müdrike Mimarisi'ndeki 41 idrak melekesini hakiki ve fiziki karşılığı olan **Reel Geometrik Cebir ($Clifford\ Algebra\ Cl_{3,0}$)**, **Reel Yatkın-Simetrik Operatörler ($J \in \mathfrak{so}(2N)$)** ve **Reel Hartley Spektral Süzgeçleri** zemininde yeniden inşa eder. Dokümanda; 41 melekenin her birinin sanal $i$ birimi içermeyen matris yapıları, Reel Schrödinger-Liouville denklemleri, Reel Komütatörler ve $p$-adic rasyonel vektör uzaylarındaki yıkıcı ve yapıcı çakışma kuralları eksiksiz biçimde vaz' edilmektedir.
\end{abstract}

\tableofcontents
\newpage

% =================================================================
\section{Giriş: Sanal Sayıların ($i = \sqrt{-1}$) İlgası ve Hakiki Reel Zemin}
% =================================================================

\subsection{Mahiyet: $i$ Biriminin Geometrik Hakikate Tercümesi}
Matematikteki $i = \sqrt{-1}$ sanal bir kurgudur. Hakiki fizikte ve kuantum mekaniğinde $i$ birimi, 2-boyutlu reel düzlemdeki $\pi/2$ (90 derece) dik dönme matrisinden veya 3-boyutlu Reel Geometrik Cebirdeki ($Cl_{3,0}$) bivektörden ($I_{12} = e_1 e_2$) başka bir şey değildir:

\begin{equation}
i \iff J = \begin{pmatrix} 0 & -1 \\ 1 & 0 \end{pmatrix} \in \mathbb{R}^{2 \times 2}, \quad J^2 = \begin{pmatrix} 0 & -1 \\ 1 & 0 \end{pmatrix} \begin{pmatrix} 0 & -1 \\ 1 & 0 \end{pmatrix} = -\mathbf{I}_2
\end{equation}

Bu ikame sayesinde, karmaşık uzay ($\mathbb{C}^N$) iki kat boyutlu **Reel Hilbert Uzayı ($\mathcal{H}_{\mathbb{R}} \cong \mathbb{R}^{2N}$)** haline getirilir. Hiçbir denklemde hayali $i$ birimi kullanılmaz.

\subsection{Gaye}
Sanal sayıların ürettiği karmaşık faz kafa karışıklığını ilga etmek; 41 melekenin tamamını **Reel Ortogonal Dönüşüm Matrisleri ($U_{\mathbb{R}} \in SO(2N)$)** ve **Reel Kliford Rotorları ($R \in Spin(N)$)** olarak çalıştırıp, sıfır sapmayla neticeye ulaşmaktır.

\subsection{Usul: Reel Schrödinger-Liouville Büküm Denklemi}
Reel durum vektörü $|\Psi(t)\rangle_{\mathbb{R}} = \begin{pmatrix} \mathbf{u}(t) \\ \mathbf{v}(t) \end{pmatrix} \in \mathbb{R}^{2N}$ için zamana bağlı evrim denklemi:

\begin{align}
\hbar \frac{d}{dt} |\Psi(t)\rangle_{\mathbb{R}} &= -J \cdot H_{\mathbb{R}} |\Psi(t)\rangle_{\mathbb{R}}, \quad H_{\mathbb{R}} = \begin{pmatrix} A & -B \\ B & A \end{pmatrix} \in \mathbb{R}^{2N \times 2N} \\
\hat{U}_{\mathbb{R}}(t) &= \exp\left( \frac{t}{\hbar} J \cdot H_{\mathbb{R}} \right) = \cos\left( \frac{H t}{\hbar} \right) \mathbf{I} + \sin\left( \frac{H t}{\hbar} \right) J \in SO(2N)
\end{align}

% =================================================================
\section{41 Melekenin Hakiki Reel Kuantum Operatör Formülasyonları}
% =================================================================

Aşağıda 41 idrak melekesinin sanal $i$ içermeyen hakiki reel matris ve operatör formülleri vaz' edilmiştir.

\subsection{1. Müşahede Kapısı ($\hat{U}_{\mathcal{O}_1}$)}
\textbf{Mahiyet:} Haricî sinyali Reel Hilbert uzayına taşıyan reel fotonik mod kipleme kapısı.
\begin{equation}
\hat{U}_{\mathcal{O}_1} = \exp\left( \theta_1(\mathbf{x}) J \cdot (\hat{x}(\mathbf{x}) + \hat{p}(\mathbf{x})) \right) \in SO(2N)
\end{equation}

\subsection{2. İllet Keşfi Kapısı ($\hat{U}_{\mathcal{O}_2}$)}
\textbf{Mahiyet:} Nedensel bağlılığı sağlayan reel yatkın-simetrik (skew-symmetric) etkileşim kapısı.
\begin{equation}
\hat{H}_{\mathcal{O}_2} = J_{\text{nedensel}} \begin{pmatrix} 0 & -1 \\ 1 & 0 \end{pmatrix} \otimes (\hat{a}_{\text{sebep}}^T \hat{a}_{\text{müsebbep}} - \hat{a}_{\text{müsebbep}}^T \hat{a}_{\text{sebep}})
\end{equation}

\subsection{3. Gaye Belirleme Kapısı ($\hat{U}_{\mathcal{O}_3}$)}
\textbf{Mahiyet:} Dalganın hedef duruma ($|x^*\rangle_{\mathbb{R}}$) odaklanmasını sağlayan reel potansiyel yansıma kapısı.
\begin{equation}
\hat{U}_{\mathcal{O}_3} = \mathbf{I} - 2 |x^*\rangle_{\mathbb{R}} \langle x^*|_{\mathbb{R}} \quad (\text{Reel Householder Yansıması})
\end{equation}

\subsection{4. Fıtrâtı İdrak Kapısı ($\hat{U}_{\mathcal{O}_4}$)}
\textbf{Mahiyet:} Fıtrî aksiyom ihlallerinde işareti değiştiren reel parite kapısı.
\begin{equation}
\hat{U}_{\mathcal{O}_4} = \mathbf{I} - 2 \hat{P}_{\text{fıtrat\_ihlali}} \quad (\text{ihlal ALT UZAYINDA işaret; genel çarpan olarak $(-1)$ ölçülemez})
\end{equation}

\subsection{5. Tecrit Kapısı ($\hat{U}_{\mathcal{O}_5}$)}
\textbf{Mahiyet:} Reel Grassmannian manifoldu ($Gr(k, n, \mathbb{R})$) üzerindeki dik izdüşüm kapısı.
\begin{equation}
\hat{U}_{\mathcal{O}_5} = \mathbf{P}_{\text{öz}} = Q Q^T, \quad Q \in \mathbb{R}^{n \times k}, \quad Q^T Q = \mathbf{I}_k
\end{equation}

\subsection{6. Tasavvur Kapısı ($\hat{U}_{\mathcal{O}_6}$)}
\textbf{Mahiyet:} Karmaşık QFT yerine, tamamen reel olan Kuantum Hartley Dönüşümü (RHT) kapısı.
\begin{align}
\text{cas}(\theta) &\coloneqq \cos(\theta) + \sin(\theta) \quad (\text{Reel Hartley Çirpması}) \\
\hat{U}_{\mathcal{O}_6} \equiv \text{RHT}_N &= \frac{1}{\sqrt{N}} \left[ \text{cas}\left( \frac{2\pi j k}{N} \right) \right]_{j,k=0}^{N-1} \in \mathbb{R}^{N \times N}
\end{align}

\subsection{7. Teemmül Kapısı ($\hat{U}_{\mathcal{O}_7}$)}
\textbf{Mahiyet:} Dalganın reel dik uzayda dönmesini sağlayan adyabatik evrim kapısı.
\begin{equation}
\hat{U}_{\mathcal{O}_7} = \mathcal{P} \exp\left( \frac{1}{\hbar} \int_0^T J \cdot H_{\text{teemmül}}(t) \, dt \right) \in SO(2N)
\end{equation}

\subsection{8. Tenakuz Bulma Kapısı ($\hat{U}_{\mathcal{O}_8}$)}
\textbf{Mahiyet:} Çelişkili durumların genliğini reel vektör uzayında zıt işarete ($-1$) katlayıp yok eden kırpma kapısı.
\begin{align}
\hat{U}_{\mathcal{O}_8} &= \mathbf{I} - 2 |\Psi_{\text{tenakuz}}\rangle_{\mathbb{R}} \langle \Psi_{\text{tenakuz}}|_{\mathbb{R}} \\
\hat{U}_{\mathcal{O}_8} |\Psi_{\text{metin}}\rangle_{\mathbb{R}} &= |\Psi_{\text{metin}}\rangle_{\mathbb{R}} - 2 \langle \Psi_{\text{tenakuz}} | \Psi_{\text{metin}} \rangle |\Psi_{\text{tenakuz}}\rangle_{\mathbb{R}} \quad (\text{yalnız tenakuz BİLEŞENİ}) \\ &\quad \text{ve } \| \hat{U}_{\mathcal{O}_8} \Psi \| = \| \Psi \| \text{ olduğundan genlik YOK OLMAZ}
\end{align}

\subsection{9. Tezat İdraki Kapısı ($\hat{U}_{\mathcal{O}_9}$)}
\textbf{Mahiyet:} Zıt iki fikri zıt eksenlere ayıran reel diyagonal kapı.
\begin{equation}
\hat{U}_{\mathcal{O}_9} = \begin{pmatrix} 1 & 0 \\ 0 & -1 \end{pmatrix}_{\text{tezat}} \otimes \mathbf{I}
\end{equation}

\subsection{10. Tahkik Kapısı ($\hat{U}_{\mathcal{O}_{10}}$)}
\textbf{Mahiyet:} Reel Matris Tekil Değer Ayrışımı (Real SVD) ile tahkik icra eden kapı.
\begin{equation}
\hat{U}_{\mathcal{O}_{10}} = \text{QSVT}_{\mathbb{R}}(A_{\mathbb{R}}, \vec{\phi}) = \sum_k P(\sigma_k) |u_k\rangle_{\mathbb{R}} \langle v_k|_{\mathbb{R}}
\end{equation}

\subsection{11. Şek-Zan-Yakîn İdraki Kapısı ($\hat{U}_{\mathcal{O}_{11}}$)}
\textbf{Mahiyet:} 3-boyutlu reel dik dönme matrisi ($SO(3)$) ile ihtimal derecelerini veren kapı.
\begin{equation}
\hat{U}_{\mathcal{O}_{11}} = \begin{pmatrix} \cos\theta & -\sin\theta & 0 \\ \sin\theta & \cos\theta & 0 \\ 0 & 0 & 1 \end{pmatrix} \in SO(3)
\end{equation}

\subsection{12. Muhakeme Kapısı ($\hat{U}_{\mathcal{O}_{12}}$)}
\textbf{Mahiyet:} Reel Kombinatoryal Hodge-Laplasyen ($\hat{\Delta}_k = \partial_{k+1} \partial_{k+1}^T + \partial_k^T \partial_k$) kapısı.
\begin{equation}
\hat{U}_{\mathcal{O}_{12}} = \exp\left( J \cdot \hat{\Delta}_k t \right) \in SO(2N)
\end{equation}

\subsection{13. İspat Kapısı ($\hat{U}_{\mathcal{O}_{13}}$)}
\textbf{Mahiyet:} Reel Kübik HoTT $Glue$ denkliği ile muhkem duruma bağlama kapısı.
\begin{equation}
\hat{U}_{\mathcal{O}_{13}} = \text{Glue}_{\mathbb{R}} \implies (\mathcal{A} \simeq_{\mathbb{R}} \mathcal{B}) \implies \hat{U}_{\mathcal{O}_{13}} |\mathcal{A}\rangle_{\mathbb{R}} = |\mathcal{B}\rangle_{\mathbb{R}}
\end{equation}

\subsection{14. Tahlil Kapısı ($\hat{U}_{\mathcal{O}_{14}}$)}
\textbf{Mahiyet:} Reel Givens dönme matrisleri ile KAN ayrışımı yapan kapı.
\begin{equation}
\hat{U}_{\mathcal{O}_{14}} = \prod_{q,p} \begin{pmatrix} \cos\phi_{q,p} & -\sin\phi_{q,p} \\ \sin\phi_{q,p} & \cos\phi_{q,p} \end{pmatrix}
\end{equation}

\subsection{15. Terkip Kapısı ($\hat{U}_{\mathcal{O}_{15}}$)}
\textbf{Mahiyet:} Reel Hartley süzgeçleri ile bileşenleri birleştiren kapı.
\begin{equation}
\hat{U}_{\mathcal{O}_{15}} = \text{RHT}^{-1} \cdot R_{\text{reel}} \cdot \text{RHT}, \quad R_{\text{reel}}[k] = R_{\text{reel}}[N-k] \ (\text{ÇİFT-simetri şart}); \text{ aksi hâlde netice bir evrişim (dolaşımlı dizey) değildir}
\end{equation}

\subsection{16. İhtimal Hesabı Kapısı ($\hat{U}_{\mathcal{O}_{16}}$)}
\textbf{Mahiyet:} Reel vektör genliklerinin karesiyle çalışan Born kuralı kapısı.
\begin{equation}
\hat{M}_{\mathcal{O}_{16}} = \sum_x \sqrt{P(x)} |x\rangle_{\mathbb{R}} \langle x|_{\mathbb{R}}, \quad P(x) = \psi(x)_{\mathbb{R}}^2 \quad (\text{ÜNİTER DEĞİL, ölçüm işleci: } \hat{M}^T \hat{M} = \mathrm{diag}(P))
\end{equation}

\subsection{17--25. Akli ve Mantıki İdrak Melekeleri Kapı Grubu ($\hat{U}_{\mathcal{O}_{17} \dots \mathcal{O}_{25}}$)}
Kıyas, Temsil, Teşbih, Temkin, Tashih, Tertip, Tafsil, Mizan, İntaç melekelerinin reel matrisleri:
\begin{align}
\hat{U}_{\text{Kıyas}} &= \begin{pmatrix} \cos\theta & \sin\theta \\ -\sin\theta & \cos\theta \end{pmatrix} \quad (\text{Reel Orantı Takas Kapısı}) \\
\hat{U}_{\text{Mizan}} &= \exp\left( J \cdot \lambda (\hat{H}_{\text{mevcut}} - \hat{H}_{\text{fıtrî}})^2 \right) \in SO(2N) \\
\hat{U}_{\text{İntaç}} &= \hat{P}_{hLevel 0} \quad (\text{Reel Kesme İntaç Kapısı})
\end{align}

\subsection{26--41. Lisan ve Belagat Kapı Grubu ($\hat{U}_{\mathcal{O}_{26} \dots \mathcal{O}_{41}}$)}
Hafıza, İrade, Kelam, Tefsir, Tevil, Fesahat, Talakat, Belagat, Sanat ve Münazara melekelerinin reel kapıları:
\begin{align}
\hat{U}_{\text{Hafıza}} &= \mathbf{I}_{\mathbb{Q}_p} \quad (\text{$p$-Adic ULTRAMETRİK Bozulmasız Saklama Kapısı}) \\
\hat{U}_{\text{Kelam}} &= U_{\text{Reel\_Fotonik}} \quad (\text{Reel Sürekli Değişkenli İntaç Kapısı}) \\
\hat{U}_{\text{Münazara}} &= \begin{pmatrix} 0 & -\mathbf{I} \\ \mathbf{I} & 0 \end{pmatrix} \quad (\text{Reel Cerh ve Teyit Matrisi})
\end{align}

% =================================================================
\section{Melekeler Arası Reel Lie Cebrail Komütatörleri}
% =================================================================

Reel ortogonal gruplardaki ($\mathfrak{so}(2N)$) Lie cebri komütatörleri, sanal $i$ çarpanı olmaksızın doğrudan yatkın-simetrik yapıyla ifade olunur:

\begin{equation}
[\hat{X}_{\mathcal{O}_i}, \hat{X}_{\mathcal{O}_j}] = \sum_{k=1}^{41} C_{ij}^k \hat{X}_{\mathcal{O}_k} \in \mathfrak{so}(2N), \quad \hat{X}_{\mathcal{O}_k} = J \cdot \hat{H}_{\mathcal{O}_k} \in \mathfrak{so}(2N), \quad \hat{U}_{\mathcal{O}_k} = e^{\hat{X}_{\mathcal{O}_k}} \in SO(2N)
\end{equation}

Bu bağıntı idrakteki mukaddemlik ve muahharık (öncelik-sonralık) sırasının kesin garantisidir.

% =================================================================
\section{Reel Vektör Uzayında Çakışma ve İlgasız İntaç}
% =================================================================

Metin dalgası $|\Psi_{\text{metin}}\rangle_{\mathbb{R}}$ ile 41 meleke dalgası arasındaki çakışma, karmaşık integraller yerine **Reel Vektör Toplamı ve Yansımaları** ile icra edilir:

\begin{equation}
|\Psi_{\text{neticeli}}\rangle_{\mathbb{R}} = \left( \prod_{k=1}^{41} \hat{U}_{\mathcal{O}_k} \right) |\Psi_{\text{metin}}\rangle_{\mathbb{R}}
\end{equation}

Çelişkili ve tenakuzlu cümle bileşenleri Householder yansımaları ($-\mathbf{I}$) sayesinde birbirini zıt vektör eklemesiyle ($v + (-v) = 0$) tamamen sıfırlar. Geriye yalnız hakiki kelam kalır.

% =================================================================
\section{Netice}
% =================================================================

Bu risalede vaz' edilen formülasyon ile uydurma sanal $i = \sqrt{-1}$ birimi kâmilen ilga edilmiştir. Nefs-i Müdrike Mimarisi'nin 41 idrak melekesi tamamıyla **Reel Geometrik Cebir ($Cl_{3,0}$)**, **Reel Ortogonal Dönüşüm Matrisleri ($SO(2N)$)** ve **Reel Hartley Spektral Süzgeçleri** zemininde yeniden vaz' edilmiş; sistem sanal sayı yanılsamasından arındırılmıştır.

\end{document}
